\begin{problemsummary}
{A page is slow, or answers \code{Unexpected Execution Error}, at HTTP 200 with
a body identical to a healthy one or an error naming a service and nothing
about why. Every service log is full of statements and none of them can be
tied to the request somebody complained about.}
{The router gives every request a trace id and sends it to each subgraph in a
W3C \code{traceparent} header with no configuration at all, and writes it on
its own log line. ASP.NET Core adopts that id in every service and attaches it
to every log scope. Nothing prints it: the console formatter leaves scopes
out, a statement log written straight to the console never sees them, the
router does not return the id to the client, and Hot Chocolate writes a resolver's exception nowhere at the log levels a
service ships with. The router's error does name the failing service
in \code{serviceName}, and names no service for a request that is merely
slow. Its caches save no service any work: a cached plan costs the same
statements, and \code{cache\_control\_policy} decides how long a whole response
may be kept from values set per service, marking a response carrying a
bearer-token column \code{public} if told to.}
{Turn on \code{response\_trace\_id} so the client receives the id in
\code{X-WG-Trace-Id}. In every service, set \code{FormatterName} to
\code{simple} beside \code{IncludeScopes} in \code{appsettings.json}, because
the option alone is ignored; print \code{Activity.Current}'s trace id on
anything that writes to the console directly; and register a diagnostic event
listener that logs resolver errors, after \code{AddApplicationService} has
registered its logger. Then a reported id finds the
slow service by counting its lines and the failing one by its exception. Leave \code{cache\_control\_policy} off unless every field a response can
carry is as cacheable as the value it would be given.}
\end{problemsummary}
